\subsection{Recombination}

While the importance of recombination in accelerating evolution has long been recognized \citep{felsenstein1976evolutionary}, a purely theoretical perspective may overlook its broader impact on the evolution of the genome itself, including its structural and regulatory organization.
recombination rates are strongly correlated with genome size \citep{tiley2015relationship}, as recombination is expected to counteract the proliferation of parasitic transposons.
Furthermore, as recombination itself is non-random and localized to hotspots determined by protein activity \citep{paul2016recombination}, selection can be modulated at finer scales.

The relationship between recombination and modularity \cite{wiser2019horizontal,kreimer2008evolution}.
For example, HGT has been associated with the expansion of the regulatory factor repertoire of bacteria \cite{price2008horizontal} and the replacement of key genes in essential metabolic pathways \cite{goyal2022horizontal}.
Recombination in sexual organisms occurs through highly regulated mechanisms that control where double-strand breaks (DSBs) are induced, resulting in non-random crossover points between parental chromosomes \cite{stapley2017recombination}.
HGT is also not limited to single or small groups of genes: a cluster of 40 photosynthetic genes has moved horizontally multiple times across alphaproteobacterial lineages as a single unit \cite{martin2018physiological, brinkmann2018horizontal}.
